理想无损同轴线内外导体半径为 \(a,b\)（\(0<a<b\)），中间均匀介质介电常数为 \(\epsilon>0\)、磁导率为 \(\mu>0\)。忽略导体电阻及介质损耗，讨论 TEM 模沿 \(+z\) 传播。定义总电压 \(V\) 为内外导体电势差，电流 \(I\) 为内导体沿 \(+z\) 电流。
(1) 从横向静电和环向磁场求单位长度电容 \(C'\)、电感 \(L'\)；(2) 从电报方程得到传播速度和特征阻抗 \(Z_0\)；(3) 当末端接纯电阻负载 \(R_L>0\) 时求电压波反射系数及在开路、短路极限的相位；(4) 通过横截面积分坡印廷矢量证明瞬时电磁功率是 \(VI\)，并说明能量在介质中流动而不必“沿导体内部传播”。
